% Phần cài đặt chi tiết của Chương 3, được \input ở cuối 03-phuong-phap.tex.
% Hai đoạn mã trích tự động bởi report/scripts/chuyen_sau/trich_ma.py; không sửa tay trong chuyen-sau/ma/.
% \par trước nhóm: nếu không, đoạn văn đứng ngay trước khối mã sẽ nhận giãn dòng đơn của khối mã.
\providecommand{\trichma}[3]{\par{\linespread{1}\selectfont
  \lstinputlisting[language=Python, caption={#2}, label={lst:cs-#3}]{chuyen-sau/ma/#1}}}

\section{Cài đặt chi tiết và các quyết định thiết kế}
\label{sec:cs-cai-dat}
Các mục trước mô tả thiết kế ở mức khái niệm. Mục này tóm tắt những quyết định cài đặt ảnh hưởng tới kết quả, gồm cả
những chỗ hành vi thực tế khác với ý đồ thiết kế. Các phép đo tương ứng nằm ở Chương~\ref{ch:thuc-nghiem}.

\subsection{Tổ chức mã nguồn}
Gói \texttt{caro} chia theo hướng phụ thuộc một chiều: \texttt{engine} không phụ thuộc gì ngoài thư viện chuẩn;
\texttt{agents} chỉ dùng \texttt{engine} (và \texttt{nn} cho cấp 5--8); \texttt{arena} chỉ dùng tác tử qua
interface; \texttt{training} dùng cả ba; \texttt{ui} đứng trên cùng. Nhờ đó một tác tử mới không đòi sửa arena, và
cấp 1--4 chạy được khi chưa cài PyTorch (bảng đăng ký chỉ import lớp khi cần). Bảng~\ref{tab:cs-thu-muc} giới thiệu
các thư mục chính, Bảng~\ref{tab:cs-tep-chinh} giới thiệu những tệp cốt lõi.
\begin{table}[H]
  \centering\small
  \caption{Các thư mục chính của dự án}
  \label{tab:cs-thu-muc}
  \begin{tabular}{lp{10cm}}
    \toprule
    Thư mục & Vai trò\\
    \midrule
    \texttt{caro/engine/} & Luật chơi: bàn cờ bitboard, trạng thái ván, khoá Zobrist, đối xứng, lõi Cython\\
    \texttt{caro/agents/} & Tác tử cấp 1--8 và bộ nối engine ngoài KataGomo, dùng chung một interface\\
    \texttt{caro/nn/} & Mã hoá thế cờ, mạng policy--value, nạp dữ liệu huấn luyện\\
    \texttt{caro/training/} & Sinh dữ liệu chuyên gia, học có giám sát, REINFORCE, vòng AlphaZero\\
    \texttt{caro/arena/} & Khai cuộc, lịch đấu, chạy giải, thống kê và đo hiệu năng\\
    \texttt{caro/ui/} & Máy chủ web và giao diện chơi, xem lại ván\\
    \midrule
    \texttt{configs/} & Cấu hình YAML của tác tử, ngân sách, huấn luyện và thí nghiệm\\
    \texttt{experiments/} & Script cho các thí nghiệm riêng\\
    \texttt{tests/} & Kiểm thử tự động\\
    \texttt{notebooks/} & Notebook phân tích và vẽ hình\\
    \texttt{runs/} & Log, cấu hình đã giải và kết quả của từng lượt chạy\\
    \texttt{report/} & Mã nguồn LaTeX của báo cáo và script sinh bảng, hình\\
    \texttt{docs/} & Đề bài, hướng dẫn và ghi chú kỹ thuật\\
    \bottomrule
  \end{tabular}
\end{table}
\begin{table}[H]
  \centering\small
  \caption{Những tệp cốt lõi của gói \texttt{caro}}
  \label{tab:cs-tep-chinh}
  \begin{tabular}{lp{9.2cm}}
    \toprule
    Tệp & Nội dung\\
    \midrule
    \texttt{engine/bitboard.py} & Kiểm tra thắng, sinh nước ứng viên, đếm cửa sổ 5 ô, lượng giá\\
    \texttt{engine/game.py} & Lớp \texttt{CaroGame}: đi, lùi, nước hợp lệ, kết quả\\
    \texttt{agents/base.py} & Interface tác tử, ngân sách mỗi nước và đồng hồ\\
    \texttt{agents/alphabeta.py} & Cấp 3: Negamax Alpha-Beta, bảng chuyển vị, sắp xếp nước\\
    \texttt{agents/mcts.py} & Cấp 4: MCTS với rollout\\
    \texttt{agents/puct.py} & Cấp 7, 8: tìm kiếm PUCT dẫn bằng mạng\\
    \texttt{nn/model.py} & Mạng \texttt{CaroResNet}\\
    \texttt{arena/tournament.py} & Chạy giải đấu và ghi log từng ván\\
    \texttt{config.py} & Ghép cấu hình và giải tham chiếu tác tử\\
    \texttt{runlog.py} & Tạo thư mục run, ghi cấu hình, seed và kết quả\\
    \texttt{hub.py} & Tải và đẩy trọng số trên Hugging Face theo phiên bản đã ghim\\
    \bottomrule
  \end{tabular}
\end{table}

\subsection{Cấu hình, seed và ngân sách}
Một tệp YAML khai báo \texttt{defaults}; các tệp cha được ghép trước theo thứ tự trái sang phải, rồi tới nội dung của
tệp và các tham số \texttt{a.b.c=giá trị} trên dòng lệnh. Từ điển được ghép đệ quy còn danh sách bị thay thế, nên
một thí nghiệm khai báo lại \texttt{agents} không giữ sót tác tử của tệp cha. Tác tử được tham chiếu bằng tên tệp kèm
phần ghi đè; thư mục run lưu bản cấu hình đã giải đầy đủ nên chạy lại không phụ thuộc các tệp cha về sau.

Seed của mỗi ván và mỗi bên là mã băm BLAKE2b của bộ (seed gốc, mã ván, màu), ổn định giữa các tiến trình. Vì vậy kết
quả ở chế độ ngân sách cố định không phụ thuộc số tiến trình song song; chế độ thời gian thì phụ thuộc tải máy.
Ngân sách được kiểm tra ngay khi tạo (từ chối giá trị âm, bằng 0 hoặc thiếu \texttt{time\_ms}), và đồng hồ trừ một
biên an toàn 10\,ms cho Alpha-Beta, 5\,ms cho MCTS và PUCT.

\subsection{Cấp 2: thứ tự luật}
Thứ tự P1--P7 chính là cách giải xung đột: luật đứng trước khớp thì các luật sau không được xét.
\begin{itemize}
  \item P3 và P4 trả nước \emph{đầu tiên} thoả điều kiện theo chỉ số ô, không so điểm giữa các nước cùng loại; khi đối
  thủ có hai đe doạ đôi độc lập, chỉ đe doạ gặp trước được chặn.
  \item P4 chỉ phòng thủ một nước, không phát hiện chuỗi tứ dài hơn, và đứng trước P5 (tạo tứ). Hệ quả được đo ở
  Bảng~\ref{tab:cs-ct-l2}.
  \item P6 và P7 chỉ khác nhau khi hoà điểm: chọn ô gần tâm, rồi ngẫu nhiên theo seed.
\end{itemize}

\subsection{Cấp 3: Alpha-Beta}
\paragraph{Một lượt quét cho cả nút.} Thông tin đe doạ được tính một lần cho mọi ứng viên theo
định lý~\ref{dl:cs-gain}. Điểm vị trí cũng vậy: đặt quân vào ô $m$ chỉ đổi các cửa sổ chứa $m$, nên
\begin{equation}
  S(M\cup\{m\},O)=S(M,O)+\sum_{h=1}^{4}\sum_{k=0}^{4}(w'_{k+1}-w'_k)\cdot\mathrm{popcount}(B_{h,k}\wedge A_{h}(m)),
\end{equation}
với $B_{h,k}$ là các cửa sổ sống hướng $h$ có $k$ quân, $A_h(m)$ là các cửa sổ hướng $h$ chứa $m$, $w'_0=w'_5=0$ và
$w'_k=w_k$ còn lại. Mỗi ứng viên chỉ tốn vài phép AND và đếm bit.

\paragraph{Trình tự trong một nút.} (1) đếm nút; (2) kiểm tra bên vừa đi đã thắng; (3) tra bảng chuyển vị; (4) lá hoặc
hết giờ thì lượng giá; (5) sinh ứng viên bán kính 2; (6) thu về nước thắng ngay, nếu không có thì các nước chặn;
(7) sắp xếp theo Bảng~\ref{tab:cs-thu-tu}; (8) duyệt, cắt tỉa, cập nhật killer và history; (9) ghi bảng nếu chưa hết
giờ. Bước (6) thu hẹp cây khi có đe doạ nhưng không tăng độ sâu, nên chuỗi tứ dài hơn độ sâu còn lại vẫn bị cắt ở lá.
\begin{table}[H]
  \centering\small
  \caption{Thang ưu tiên sắp xếp nước của cấp 3}
  \label{tab:cs-thu-tu}
  \begin{tabular}{rl}
    \toprule
    Điểm & Lớp nước\\
    \midrule
    10.000.000 & Nước lưu trong bảng chuyển vị cho đúng thế này\\
    9.500.000 & Thắng ngay\\
    9.000.000 & Tạo từ hai ô thắng trở lên (đòn kép)\\
    8.000.000 & Tạo tứ, hoặc đang có tứ\\
    7.000.000 & Chiếm ô mà đối thủ dùng để tạo đòn kép\\
    6.000.000 / 5.000.000 & Hai nước killer của ply\\
    1.000.000 $+h(m)$ & Lịch sử cắt tỉa, $h$ cộng $d^2$ mỗi lần cắt\\
    dưới 1.000.000 & Điểm vị trí sau khi đi\\
    \bottomrule
  \end{tabular}
\end{table}

\paragraph{Gốc tìm kiếm và lặp sâu dần.}
\begin{itemize}
  \item Danh sách nước gốc không qua bộ lọc nước bắt buộc. Ở độ sâu 1, một nước không chặn ô thắng của đối thủ vẫn có
  thể được chấm cao; các giải đấu dùng độ sâu từ 4 nên không gặp (xem Bảng~\ref{tab:cs-ct-l3}).
  \item Nước tốt nhất của độ sâu trước được đưa lên đầu danh sách, nhưng hàm gốc lại sắp xếp danh sách đó nên thứ tự
  này bị ghi đè; hệ quả đo ở Mục~\ref{sec:cs-goc}.
  \item Hết giờ giữa chừng thì bỏ kết quả của độ sâu dở dang; tìm được đường thắng ép thì dừng ngay.
\end{itemize}

\paragraph{Bảng chuyển vị và giải hoà.} Bảng có kích thước cố định (64\,MB, ước lượng 200 byte mỗi mục), khi đầy thì
đẩy mục được chèn sớm nhất. Khi hoà điểm, bước sắp xếp đặt ô chỉ số lớn trước, còn hàm sinh nhãn
\texttt{scored\_moves} giữ ô chỉ số nhỏ trước; chi tiết thứ hai gây ra lỗi nhãn ở Mục~\ref{sec:cs-du-lieu}.

\paragraph{Lõi Cython.} Bản native dịch từng dòng sang \texttt{unsigned \_\_int128}, giữ cả thứ tự duyệt, cách giải
hoà và chính sách thay thế của bảng. Hai backend vì vậy cho cùng điểm, biến chính và số nút
(\texttt{tests/test\_native.py}), chỉ khác tốc độ (Bảng~\ref{tab:cs-tk-backend-ab}).

\subsection{Cấp 4: MCTS}
Mỗi mô phỏng mở rộng một con chọn ngẫu nhiên trong các nước chưa thử, rồi chơi rollout đến hết ván. Chế độ
\texttt{random} không xét thắng ngay; chế độ \texttt{light} thắng ngay rồi chặn trước khi đi ngẫu nhiên. $W$ của mỗi
nút theo góc nhìn bên đã đi vào nút, nên công thức chọn dùng $W/N$ trực tiếp. Với khoảng 80 ứng viên ở gốc và ngân
sách 400 mô phỏng, mỗi con chỉ có khoảng 5 mẫu, rất xa mức mà định lý~\ref{dl:cs-hoeffding} đòi hỏi. Nước được chọn là
con có $N$ lớn nhất. Tuỳ chọn \texttt{root\_tactics} (A8) xử lý ở gốc việc thắng ngay, chặn và đòn kép trước khi tìm;
đó là heuristic chứ không phải bộ giải, vì một tứ phản công cũng là phòng thủ.

\subsection{Mạng và PUCT (cấp 5--8)}
Đầu vào gồm 4 kênh: quân bên đi, quân đối thủ, hằng 1 nếu X đi, và một kênh luôn bằng 0 (không mang thông tin nhưng
vẫn chiếm trọng số ở lớp vào). Thân mạng gồm sáu khối dư 64 kênh; ô đã có quân bị gán $-\infty$ trước softmax, đầu
value kết thúc bằng $\tanh$. Phân bố tham số cho trong Bảng~\ref{tab:cs-mang-thamso}.
\begin{table}[H]
  \centering
  \footnotesize
  \caption{Phân bố tham số của CaroResNet (64 kênh, 6 khối)}
  \label{tab:cs-mang-thamso}
  \begin{tabular}{lrr}
    \toprule
    Khối & Tham số & \% \\
    \midrule
    \texttt{conv\_in} & 2.304 & 0,5 \\
    \texttt{bn\_in} & 128 & 0,0 \\
    \texttt{res\_blocks.0} & 73.984 & 15,4 \\
    \texttt{res\_blocks.1} & 73.984 & 15,4 \\
    \texttt{res\_blocks.2} & 73.984 & 15,4 \\
    \texttt{res\_blocks.3} & 73.984 & 15,4 \\
    \texttt{res\_blocks.4} & 73.984 & 15,4 \\
    \texttt{res\_blocks.5} & 73.984 & 15,4 \\
    \texttt{policy\_conv} & 128 & 0,0 \\
    \texttt{policy\_bn} & 4 & 0,0 \\
    \texttt{policy\_fc} & 29.403 & 6,1 \\
    \texttt{value\_conv} & 64 & 0,0 \\
    \texttt{value\_bn} & 2 & 0,0 \\
    \texttt{value\_fc1} & 3.904 & 0,8 \\
    \texttt{value\_fc2} & 33 & 0,0 \\
    \midrule
    Tổng & 479.874 & 100,0 \\
    \bottomrule
  \end{tabular}
\end{table}

PUCT lưu $W$ theo góc nhìn bên đang đi tại nút, ngược với cấp 4, nên giá trị của con nhìn từ cha là $-W/n$:
\trichma{puct-chon.py}{Quy tắc chọn con của PUCT, có virtual loss}{puct-chon}
Tìm kiếm viết dạng generator: mỗi khi cần mạng, nó trả ra một lô thế cờ và nhận lại kết quả, nhờ đó cùng một thuật
toán chạy được cho một ván hoặc ghép suy luận của nhiều ván tự đấu. Mỗi lượt gom lá đi xuống tới nút chưa mở, cộng
virtual loss để các lá khác nhau, suy luận một lần cho cả lô rồi lan ngược; lá trùng trong lô chỉ mở rộng một lần. Bản
Cython chỉ chuyển việc chọn và tạo con, biên dịch với \texttt{-ffp-contract=off} để điểm trùng từng bit với Python.

\subsection{Huấn luyện}
\paragraph{Dữ liệu chuyên gia.} Ở mỗi thế, Alpha-Beta chấm các ứng viên, giữ top-$k$, lấy softmax theo nhiệt độ làm
nhãn policy rồi lấy mẫu nước đi từ chính phân bố đó. Nhãn value gán sau khi ván kết thúc. Thế trùng nhau qua tám phép
đối xứng được loại.

\paragraph{Học có giám sát.} Nhãn value được đổi sang góc nhìn bên đi, và tập được chia theo \emph{ván} để các thế gần
giống nhau không rơi vào hai tập:
\trichma{sl-chia.py}{Đổi nhãn value và chia tập theo ván}{sl-chia}
Khi huấn luyện, mỗi mẫu train được áp ngẫu nhiên một phép đối xứng cho cả trạng thái và nhãn. Hàm mất mát là entropy
chéo với nhãn mềm cộng MSE của value, tối ưu bằng Adam; mạng có loss validation thấp nhất được lưu.

\paragraph{REINFORCE.} Mỗi vòng chơi 128 ván song song: một nửa tự đấu với bản mới nhất, nửa còn lại gặp một trong năm
bản cũ gần nhất để giữ đa dạng đối thủ. Lợi thế $z-v$ không truyền gradient vào value; gradient được cộng dồn theo lô
con và cắt chuẩn ở 1.

\paragraph{Vòng AlphaZero.} Tự đấu dùng nhiễu Dirichlet ($\alpha=0{,}3$, $\varepsilon=0{,}25$) ở gốc và lấy mẫu theo lượt
thăm trong 8 nước đầu \cite{silver2018alphazero}. Huấn luyện lấy mẫu đều theo thế cờ trong cửa sổ replay. Mạng mới phải
đạt $0{,}55$ trước mạng đương nhiệm; nếu không, trọng số và trạng thái Adam được khôi phục nhưng các ván tự đấu vẫn
được giữ.

\subsection{Arena, giao diện và engine ngoài}
Lịch đấu tất định: hai ván liên tiếp dùng chung khai cuộc và đổi màu. Thời gian mỗi nước đo ở arena, ngoài tác tử. Bộ
chạy giải tạo mới tác tử cho từng ván; trường bộ nhớ nó ghi chỉ là mức tăng đỉnh RSS nên báo cáo đo RAM riêng. Tệp
\texttt{stats.json} vẫn ghi Elo tuần tự để xem nhanh; mọi bảng Elo trong báo cáo dùng Bradley-Terry
(Bảng~\ref{tab:cs-tk-elo}).

Giao diện web dùng đúng tác tử và cấu hình của arena; nước của máy chạy trong thread pool, nên ngân sách bị chặn trên.
Bảng~\ref{tab:cs-api} liệt kê các điểm cuối.
\begin{table}[H]
  \centering\small
  \caption{API của máy chủ giao diện}
  \label{tab:cs-api}
  \begin{tabular}{p{6.4cm}p{8cm}}
    \toprule
    Điểm cuối & Chức năng\\
    \midrule
    \texttt{GET /agents}, \texttt{/budgets}, \texttt{/openings} & Danh sách tác tử, chế độ ngân sách, khai cuộc\\
    \texttt{POST /games} & Tạo ván người--máy hoặc máy--máy\\
    \texttt{GET /games/\{id\}} & Trạng thái ván và thống kê tìm kiếm gần nhất\\
    \texttt{POST /games/\{id\}/move} (\texttt{ai-move}, \texttt{undo}) & Người đi, máy đi, lùi nước\\
    \texttt{GET /replays}, \texttt{/replays/\{id\}} & Liệt kê và phát lại ván từ log\\
    \bottomrule
  \end{tabular}
\end{table}
Tác tử KataGomo giữ một tiến trình GTP và đồng bộ lại cả ván trước mỗi nước; engine phải dùng nhánh đã vá để đúng luật
năm quân freestyle (\repoFile{docs/KATAGOMO.md}) \cite{katagomo}.

\subsection{Kiểm thử}
Bộ test trong \repoFile{tests/} kiểm ba kiểu tính chất: đối chiếu engine với bản duyệt mảng ngây thơ; đối chiếu hai cài
đặt của cùng một thuật toán (Python và Cython, mã hoá đặc trưng đơn và theo lô); và các bất biến như đếm cây trò chơi
(perft), đối xứng, lùi nước khôi phục khoá. Ngoài ra có test cho cấu hình, arena, thống kê, huấn luyện, giao diện web
và bộ giải VCF dùng trong phân tích.
